\documentclass[12pt, a4paper]{article}
\usepackage[utf8]{inputenc}
\usepackage[T1]{fontenc}
\usepackage{amsmath, amssymb, amsthm}
\usepackage{geometry}
\geometry{margin=2.5cm}

\newtheorem{lema}{Lema}
\newtheorem{proposicao}{Proposição}
\newtheorem{observacao}{Observação}

\title{Uma propriedade de soma constante na grade $3\times N$ e as conjecturas de Goldbach}
\author{Tiago Bandeira}
\date{Maio de 2026}

\begin{document}
	
	\maketitle
	
	\begin{abstract}
		A grade $G_{3,N}$, formada pelos primeiros $3N$ números ímpares organizados em três linhas, possui uma propriedade aritmética notável: em cada janela $3\times 3$ as três configurações canônicas (diagonal principal, diagonal secundária e coluna central) têm a mesma soma $S = 3(2N+2i+1)$. Quando tais configurações são totalmente primas, os pares de elementos fornecem representações da Conjectura de Goldbach (forte) e a soma dos três elementos fornece uma representação da Conjectura de Goldbach fraca. Este artigo explora essa propriedade geométrico-aditiva de forma elementar e autocontida.
	\end{abstract}
	
	\section{Introdução}
	
	A grade $G_{3,N}$ é construída organizando os números ímpares consecutivos em três linhas de $N$ colunas. Em trabalhos anteriores~\cite{bandeira2026}, essa grade foi usada para estudar a distribuição de primos em configurações rígidas. Durante essas investigações, observou-se uma identidade algébrica simples, porém surpreendente: em cada submatriz $3\times 3$ (denominada janela $W_i$), as três configurações que desempenham papel central – a diagonal principal, a diagonal secundária e a coluna central – possuem sempre a mesma soma.
	
	Este artigo apresenta essa propriedade de forma independente, demonstra-a rigorosamente e mostra como ela fornece uma \textit{testemunha geométrica unificada} das duas conjecturas de Goldbach:
	\begin{itemize}
		\item a \textbf{conjectura forte} (binária): todo número par $>2$ é soma de dois primos;
		\item a \textbf{conjectura fraca} (ternária): todo número ímpar $>5$ é soma de três primos.
	\end{itemize}
	Enquanto a conjectura fraca já foi demonstrada por Helfgott~\cite{helfgott2013}, a conexão estrutural aqui apresentada permanece interessante do ponto de vista didático e heurístico.
	
	\section{Definição da grade $G_{3,N}$ e das janelas}
	
	\subsection{Grade $G_{3,N}$}
	Seja $N\ge 3$ um inteiro. A grade $G_{3,N}$ é uma matriz com $3$ linhas e $N$ colunas, cujas células recebem os números ímpares positivos em ordem linha a linha. A célula na linha $r\in\{1,2,3\}$ e coluna $c\in\{1,\dots,N\}$ é definida por
	\[
	f(r,c)=2\bigl[(r-1)N + c\bigr] - 1.
	\]
	Assim, as três linhas contêm respectivamente:
	\begin{align*}
		\text{Linha 1:}&\quad 1,\;3,\;\dots,\;2N-1,\\
		\text{Linha 2:}&\quad 2N+1,\;2N+3,\;\dots,\;4N-1,\\
		\text{Linha 3:}&\quad 4N+1,\;4N+3,\;\dots,\;6N-1.
	\end{align*}
	
	\subsection{Janelas $3\times 3$}
	Para $i=1,\dots,N-2$, a \textbf{janela} $W_i$ é a submatriz de $G_{3,N}$ formada pelas colunas $i$, $i+1$ e $i+2$. Dentro de $W_i$ destacamos três configurações canônicas (ver Figura~\ref{fig:janela}):
	\begin{itemize}
		\item $\mathbf{D_1}$ (diagonal principal): células $(1,i)$, $(2,i+1)$, $(3,i+2)$;
		\item $\mathbf{D_2}$ (diagonal secundária): células $(3,i)$, $(2,i+1)$, $(1,i+2)$;
		\item $\mathbf{C_v}$ (coluna central): células $(1,i+1)$, $(2,i+1)$, $(3,i+1)$.
	\end{itemize}
	
	\begin{figure}[h]
		\centering
		\fbox{\begin{minipage}{0.5\textwidth}
				\centering
				\small
				\begin{tabular}{ccc}
					$f(1,i)$ & $f(1,i+1)$ & $f(1,i+2)$\\
					$f(2,i)$ & $f(2,i+1)$ & $f(2,i+2)$\\
					$f(3,i)$ & $f(3,i+1)$ & $f(3,i+2)$
				\end{tabular}
				\\[2mm]
				\emph{Diagonal D$_1$: $(1,i),(2,i+1),(3,i+2)$;\quad D$_2$: $(3,i),(2,i+1),(1,i+2)$;\quad C$_v$: $(1,i+1),(2,i+1),(3,i+1)$.}
		\end{minipage}}
		\caption{Janela $W_i$ e as três configurações.}
		\label{fig:janela}
	\end{figure}
	
	\section{Lema da soma constante}
	
	\begin{lema}
		Para qualquer janela $W_i$ ($1\le i\le N-2$) da grade $G_{3,N}$, as três configurações $D_1$, $D_2$ e $C_v$ têm a mesma soma, dada por
		\[
		\boxed{\;D_1 = D_2 = C_v = 3(2N+2i+1)\; }.
		\]
	\end{lema}
	
	\begin{proof}
		Usando a definição de $f(r,c)$:
		\begin{align*}
			f(1,i) &= 2i-1,\\
			f(2,i+1) &= 2\bigl[N+(i+1)\bigr]-1 = 2N+2i+1,\\
			f(3,i+2) &= 2\bigl[2N+(i+2)\bigr]-1 = 4N+2i+3.
		\end{align*}
		Somando os três elementos de $D_1$:
		\[
		D_1 = (2i-1)+(2N+2i+1)+(4N+2i+3)=6N+6i+3 = 3(2N+2i+1).
		\]
		
		Para $D_2$:
		\begin{align*}
			f(3,i) &= 2[2N+i]-1 = 4N+2i-1,\\
			f(2,i+1) &= 2N+2i+1,\\
			f(1,i+2) &= 2(i+2)-1 = 2i+3.
		\end{align*}
		Somando:
		\[
		D_2 = (4N+2i-1)+(2N+2i+1)+(2i+3)=6N+6i+3.
		\]
		
		Para $C_v$:
		\begin{align*}
			f(1,i+1) &= 2(i+1)-1 = 2i+1,\\
			f(2,i+1) &= 2N+2i+1,\\
			f(3,i+1) &= 2[2N+(i+1)]-1 = 4N+2i+1.
		\end{align*}
		Somando:
		\[
		C_v = (2i+1)+(2N+2i+1)+(4N+2i+1)=6N+6i+3.
		\]
		
		Portanto, as três somas coincidem. 
	\end{proof}
	
	\section{Conexão com as conjecturas de Goldbach}
	
	\subsection{Conjectura forte (binária)}
	Se uma configuração (por exemplo $D_1$) for \textbf{totalmente prima} – ou seja, $f(1,i), f(2,i+1), f(3,i+2)$ são todos primos – então os três pares de elementos fornecem números pares:
	\begin{align*}
		p_1 &= f(1,i)+f(2,i+1) = 2(N+2i),\\
		p_2 &= f(1,i)+f(3,i+2) = 2(2N+2i+1),\\
		p_3 &= f(2,i+1)+f(3,i+2) = 2(3N+2i+2).
	\end{align*}
	Cada uma dessas somas é um número par maior que $2$ (para $N\ge 3$). Assim, um único trio primo na grade produz \textbf{três representações de Goldbach binárias} (três pares de primos que somam pares distintos). O mesmo raciocínio aplica-se a $D_2$ e $C_v$: cada configuração totalmente prima gera três representações binárias, mas com valores de pares diferentes, pois os elementos que compõem cada configuração são distintos entre si.
	
	\subsection{Conjectura fraca (ternária)}
	O Lema da soma constante mostra que a soma dos três elementos da configuração é sempre $3(2N+2i+1)$. Se os três elementos forem primos, então
	\[
	\text{primo}_1 + \text{primo}_2 + \text{primo}_3 = 3(2N+2i+1),
	\]
	que é um número ímpar (pois $2N+2i+1 = 2(N+i)+1$ é ímpar, e o produto de $3$ por um número ímpar é ímpar) e maior que $5$ quando $N\ge 3$. Logo, cada trio primo na grade também fornece uma \textbf{representação da Goldbach ternária} para aquele ímpar.
	
	\begin{observacao}
		Enquanto a Goldbach fraca já foi provada por Helfgott~\cite{helfgott2013}, a estrutura geométrica aqui apresentada tem valor didático e estrutural: ela exibe, no mesmo objeto, testemunhas das duas conjecturas. Além disso, o lema independe da primalidade – a igualdade das somas é um fato algébrico da grade.
	\end{observacao}
	
	\section{Exemplo numérico}
	
	Tomemos $N=5$ e $i=2$. A grade $G_{3,5}$ é
	\[
	\begin{pmatrix}
		1 & 3 & 5 & 7 & 9\\
		11 & 13 & 15 & 17 & 19\\
		21 & 23 & 25 & 27 & 29
	\end{pmatrix}.
	\]
	A janela $W_2$ (colunas 2,3,4) é
	\[
	\begin{pmatrix}
		3 & 5 & 7\\
		13 & 15 & 17\\
		23 & 25 & 27
	\end{pmatrix}.
	\]
	Calculamos as três configurações:
	\begin{align*}
		D_1 &= 3+15+27 = 45,\\
		D_2 &= 23+15+7 = 45,\\
		C_v &= 5+15+25 = 45.
	\end{align*}
	De fato, $3(2N+2i+1)=3(10+4+1)=45$. Nenhuma delas é totalmente prima (15 e 25 não são primos), mas a identidade permanece.
	
	Para um exemplo efetivo com primos, tome $N=20$ e $i=3$. Os elementos da diagonal $D_1$ são:
	\begin{align*}
		f(1,3) &= 2\cdot 3 - 1 = 5,\\
		f(2,4) &= 2[20+4]-1 = 47,\\
		f(3,5) &= 2[40+5]-1 = 89.
	\end{align*}
	Os três são primos, e a soma constante vale
	\[
	D_1 = 5 + 47 + 89 = 141 = 3(2\cdot 20 + 2\cdot 3 + 1) = 3\cdot 47.\;\checkmark
	\]
	\textbf{Representações binárias (Goldbach forte):}
	\[
	52 = 5+47,\qquad 94 = 5+89,\qquad 136 = 47+89.
	\]
	\textbf{Representação ternária (Goldbach fraca):}
	\[
	141 = 5 + 47 + 89,\quad \text{com } 141 \text{ ímpar} > 5.\;\checkmark
	\]
	Um único trio primo na grade testemunha simultaneamente três instâncias da conjectura forte e uma da conjectura fraca.
	
	\section{Conclusão}
	
	Apresentamos uma propriedade aritmética elementar da grade $3\times N$: a diagonal principal, a diagonal secundária e a coluna central de qualquer janela $3\times 3$ têm a mesma soma $3(2N+2i+1)$. Essa igualdade transforma cada trio de primos posicionado em uma dessas configurações em uma testemunha simultânea das conjecturas de Goldbach forte (através dos pares) e fraca (através da soma tripla). A observação é independente de modelos probabilísticos e constitui um belo exemplo de como uma construção geométrica simples pode unificar dois problemas clássicos da teoria dos números.
	
	\begin{thebibliography}{9}
		\bibitem{bandeira2026}
		T. Bandeira, ``Uma Abordagem Unificada para a Conjectura de Goldbach: Estrutura Combinatória, Saturação Geométrica e o Elo entre Trios e Pares Primos'', preprint, 2026.
		
		\bibitem{bandeira2026b}
		T. Bandeira, ``Saturação Geométrica e Inevitabilidade Estrutural: Uma Abordagem de Malhas $3\times N$ para a Conjectura de Goldbach'', preprint, 2026.
		
		\bibitem{helfgott2013}
		H. A. Helfgott, ``Major arcs for Goldbach's theorem'', arXiv:1305.2897, 2013.
	\end{thebibliography}
	
\end{document}